%% ma: L12-B08-0196
%% lop: 12
%% chuong: 2
%% bai: 8
%% dang: 12.8.16
%% loai: TN4
%% muc: VD
%% dap_an: "A"
%% nguon: "CĐ2.6.pdf giáo viên gửi 27/10/2026 (M07, Lê Minh Khôi – Nguyễn Anh Khôi, Chương II Chủ đề 6 Ứng dụng vào các bài toán thực tế), Câu 32"
%% kien_thuc: "Bài 7 (tọa độ điểm khi gắn hệ trục), lớp 11 (góc nhị diện, hình chóp có các cạnh bên bằng nhau)"
%% trang_thai: chinh-thuc
%% ghi_chu: "Đề gốc ghi chiều dài 1200 cm, chiều rộng 900 cm; đã ghi rõ MN = 1200 cm, NK = 900 cm theo hình gốc"
\begin{ex}
Một ngôi nhà gồm hai phần. Phần thân nhà dạng hình hộp chữ nhật $ABCD.OMNK$ có chiều dài $MN=1200$ cm, chiều rộng $NK=900$ cm, chiều cao $450$ cm. Phần mái nhà dạng hình chóp $S.ABCD$ có các cạnh bên bằng nhau và $(SBC)$ tạo với mặt đáy góc $\alpha$ với $\tan\alpha=\dfrac13$. Chọn hệ trục tọa độ $Oxyz$ (đơn vị dài trên các trục là cen-ti-mét) sao cho $M$ thuộc tia $Ox$, $K$ thuộc tia $Oy$, $A$ thuộc tia $Oz$ (như hình vẽ). Biết tọa độ của điểm $S$ là $S(a;b;c)$. Tính giá trị của biểu thức $P=a+b+c$.
\begin{center}
\begin{tikzpicture}[scale=.42,x={(-.55cm,-.45cm)},y={(1cm,0cm)},z={(0cm,1cm)}]
  \path (0,0,0) coordinate (O) (9,0,0) coordinate (M) (9,12,0) coordinate (N) (0,12,0) coordinate (K) (0,0,4.5) coordinate (A) (9,0,4.5) coordinate (B) (9,12,4.5) coordinate (C) (0,12,4.5) coordinate (D) (4.5,6,6) coordinate (S);
  \draw (M)--(N)--(C)--(B)--cycle (N)--(K)--(D)--(C) (B)--(S)--(C) (S)--(D);
  \draw[khuat] (O)--(M) (O)--(K) (O)--(A)--(B) (A)--(D) (A)--(S);
  \draw[truc] (K)--(0,15,0) node[right]{$y$}; \draw[truc] (A)--(0,0,8) node[left]{$z$}; \draw[truc] (M)--(12.5,0,0) node[left]{$x$};
  \path (O) node[above left]{$O$} (M) node[left]{$M$} (N) node[below]{$N$} (K) node[below right]{$K$} (A) node[above left]{$A$} (B) node[left]{$B$} (C) node[below right]{$C$} (D) node[right]{$D$} (S) node[above]{$S$}
        ($(M)!.5!(N)$) node[below]{$1200$ cm} ($(N)!.5!(K)$) node[below right]{$900$ cm} ($(K)!.5!(D)$) node[right]{$450$ cm};
\end{tikzpicture}
\end{center}
\choice{\True $P=1650$}{$P=1500$}{$P=1800$}{$P=1950$}
\loigiai{$M(900;0;0)$, $K(0;1200;0)$, $A(0;0;450)$; $B(900;0;450)$, $C(900;1200;450)$. Các cạnh bên bằng nhau nên $S$ ở trên tâm $J(450;600;450)$ của $ABCD$.

$d(J;BC)=450$, $SJ=450\tan\alpha=150$. $S(450;600;600)$, $P=1650$.}
\end{ex}
